\section{An explicit cubic compositum with two excess left branches}
\label{sec:analytic-cubic-left-imbalance}

Let $N$ be the quaternionic field of NN-A018 and NN-A033.
Put $p=751$, and let $\vartheta$ be the primitive cubic
Dirichlet character modulo $p$ specified by
\[
  \vartheta(3^j)=e^{2\pi ij/3},\qquad
  \vartheta(n)=0\quad(p\mid n).
\]
Let $E$ be the fixed field of $\ker\vartheta$ in
$\mathbb Q(\zeta_{751})$, and define $K=NE$.
This explicitly specifies the field; no search over number
fields is part of the statement.

\begin{theorem}[NN-A034: an unconditional arithmetic left imbalance]
\label{thm:analytic-cubic-left-imbalance}
The extension $K/\mathbb Q$ is totally real, has Galois group
$Q_8\times C_3$, degree $24$, and absolute discriminant
\[
                   D_K=205^{18}751^{16}.
\]
There is an integer $k\ge0$ such that
\[
 \operatorname{ord}_{s=1/2}\zeta_K(s)=m=2+4k,
 \qquad \frac{\zeta_K^{(m)}(1/2)}{m!}<0,
 \qquad C_K>0.
\]
For all sufficiently small positive $t$, the shifted ideal-norm
deformation of NN-A004 has exactly $2+2k$ central branches on
the left of $\Re s=1/2$ and exactly $2k$ on the right. Thus it
has exactly two excess left branches. Its two initially vertical
branches both bend strictly left at the second Puiseux order.
All these conclusions are unconditional; neither A nor M is
used. The integer $k$, a central order of a complex Artin
twist, is not determined by this theorem.
\end{theorem}

\begin{lemma}[Four certified nonzero central values]
\label{lem:analytic-cubic-four-values}
Write $\chi_d(n)=(d/n)$ for $d=1,5,41,205$, with
$\chi_1=1$. Each $\eta_d=\chi_d\vartheta$ is primitive and
even, of conductor $q_d=751d$. Its central value lies in the
following open rectangle. Every decimal endpoint in this table
is an exact rational number.
\[
\begin{array}{c|r|c|c}
d&q_d&\Re L(1/2,\eta_d)&\Im L(1/2,\eta_d)\\\hline
1&751&(-0.0265980,-0.0265979)&(0.2188449,0.2188450)\\
5&3755&(1.2071682,1.2071683)&(-0.5141633,-0.5141632)\\
41&30791&(0.3418888,0.3418889)&(0.0415524,0.0415525)\\
205&153955&(0.8806422,0.8806423)&(2.0675988,2.0675989)
\end{array}
\]
In particular all four values are nonzero.
\end{lemma}

\begin{proof}
Primality of $751$ follows by trial division by the primes at
most $\sqrt{751}<28$. Since $750=2\cdot3\cdot5^3$ and
\[
 3^{375}\equiv750,\qquad3^{250}\equiv678,\qquad
 3^{150}\equiv80\pmod{751},
\]
the residue class of $3$ is primitive. This defines the character
exactly by its finite discrete-log table. Its parity is even
because $\vartheta(-1)=e^{2\pi i375/3}=1$. The quadratic
characters have positive fundamental discriminants $d$, and
their conductors are coprime to $751$; the product characters
are consequently primitive of conductor $q_d$ and even.

Here is the exact analytic formula and rigorous tail used for
the finite check. For any primitive even character $\eta$ of
conductor $q$, put
\[
 \epsilon_\eta=\frac1{\sqrt q}
          \sum_{a=1}^{q-1}\eta(a)e^{2\pi ia/q},\qquad
 w_q(n)=\frac{\Gamma(1/4,\pi n^2/q)}{\Gamma(1/4)\sqrt n}.
\]
The primitive Gauss sum has modulus $\sqrt q$, so
$|\epsilon_\eta|=1$. The completed functional equation has
this sign in the convention
$\Lambda(s,\eta)=(q/\pi)^{s/2}\Gamma(s/2)L(s,\eta)$;
compare DLMF~25.15.5--25.15.6. To derive the central formula,
apply Poisson summation in the residue classes to
$\Theta_\eta(x)=\sum_{n\in\mathbb Z}\eta(n)e^{-\pi n^2x/q}$.
The finite Fourier transform is
\[
 \sum_{a\bmod q}\eta(a)e^{2\pi ian/q}
              =\bar\eta(n)\tau(\eta),
\]
so
$\Theta_\eta(x)=\epsilon_\eta x^{-1/2}
\Theta_{\bar\eta}(1/x)$. Splitting its Mellin integral at $1$
gives the absolutely convergent identity
\begin{equation}\label{eq:analytic-cubic-central-sum}
 L(1/2,\eta)=\sum_{n\ge1}
            \bigl(\eta(n)+\epsilon_\eta\bar\eta(n)\bigr)w_q(n).
\end{equation}
The change of variable $u=\pi n^2x/q$ evaluates each integral
as the upper incomplete gamma function displayed in $w_q$;
dividing by $(q/\pi)^{1/4}\Gamma(1/4)$ leaves precisely
$n^{-1/2}$. This verifies every scale and phase in
\eqref{eq:analytic-cubic-central-sum}.

For $x>0$,
\[
 \Gamma(1/4,x)=\int_x^\infty u^{-3/4}e^{-u}\,du
                  \le x^{-3/4}e^{-x}.
\]
Also
$\Gamma(1/4)>\int_0^1u^{-3/4}(1-u)\,du=16/5>1$.
Consequently the omitted tail after an integer $B\ge1$ is
bounded by
\begin{equation}\label{eq:analytic-cubic-tail}
 \begin{split}
 \left|\sum_{n>B}
       (\eta(n)+\epsilon_\eta\bar\eta(n))w_q(n)\right|
 &\le 2(q/\pi)^{3/4}
       \int_B^\infty x^{-2}e^{-\pi x^2/q}\,dx\\
 &\le\frac{(q/\pi)^{7/4}}{B^3}
                         e^{-\pi B^2/q}.
 \end{split}
\end{equation}
The first comparison uses the decreasing positive summand;
the second uses $x^{-2}\le x/B^3$ for $x\ge B$ and integrates
the resulting Gaussian derivative exactly.

Only the Gauss sum of conductor $751$ is needed. The Chinese remainder
theorem, writing a residue modulo $751d$ as $751a+db$, gives
\[
 \tau(\chi_d\vartheta)
   =\chi_d(751)\vartheta(d)\tau(\chi_d)\tau(\vartheta),
 \qquad
 \epsilon_{\eta_d}
   =\chi_d(751)\vartheta(d)\epsilon_\vartheta.
\]
Here $\tau(\chi_d)=\sqrt d$, with the positive quadratic Gauss
sign verified in Gross, section~1, equations~(1.3)--(1.6),
as used in NN-A032; for $d=1$ this is the trivial identity.
This exact formula avoids any ambiguity about conjugating a
root number.

The saved script \path{work/analytic/check_cubic_compositum.py}
uses these formulas with FLINT/Arb complex balls at $60$
decimal digits. It takes
\[
 B=5(\lfloor\sqrt q\rfloor+1),\qquad
 B\in\{140,310,880,1965\},
\]
computes the full finite sum in
\eqref{eq:analytic-cubic-central-sum}, and adds the upper bound
\eqref{eq:analytic-cubic-tail} to the error in both real and
imaginary components. The respective tail bounds are strictly
less than
\[
 1.31\cdot10^{-38},\quad9.86\cdot10^{-38},\quad
 6.87\cdot10^{-37},\quad1.29\cdot10^{-36}.
\]
It checks all eight real and imaginary interval inclusions in
the table, as well as nonvanishing, and writes the complete
balls and tail bounds to
\path{work/analytic/cubic_compositum_certificate.json}.
The stated rectangles are certified enclosures, not numerical
approximations substituted for a proof.
\end{proof}

\begin{proof}[Proof of Theorem~\ref{thm:analytic-cubic-left-imbalance}]
The cyclic field $E$ exists by the cyclotomic Galois
correspondence. It has degree three and is real because
$\vartheta(-1)=1$. The only ramified prime is $751$, where the
ramification is total and tame of index three, so
$D_E=751^2$. This also follows directly from the
conductor--discriminant formula for its two nontrivial cubic
characters. The degrees $8$ and $3$ are coprime, so
$N\cap E=\mathbb Q$. Both fields are Galois, which gives
$\operatorname{Gal}(K/\mathbb Q)=Q_8\times C_3$ and degree
$24$. The discriminants $205^6$ and $751^2$ are coprime.
The integral tensor-product discriminant formula therefore
gives
\[
 D_K=D_N^3D_E^8=205^{18}751^{16}.
\]
A precise primary exposition is B.~Conrad,
\emph{Discriminant of composite fields}, Math.~154 handout,
p.~1: its displayed formula and trace-tensor determinant
calculation apply with exactly these degree and coprimality
hypotheses. Equivalently, the tame inertia indices $4,4,3$ at
$5,41,751$ give exponents $24(1-1/4)=18$ and
$24(1-1/3)=16$ directly.

Let $\psi$ denote the unique degree-two irreducible character
of $Q_8$, inflated where appropriate. The irreducible
characters of the product group are
\[
 \chi_d\vartheta^j\quad(d=1,5,41,205;\ j=0,1,2),
 \qquad
 \psi\vartheta^j\quad(j=0,1,2).
\]
The former have degree one and the latter degree two.
Decomposing the regular representation, Artin formalism gives
the exact factorization, including ramified Euler factors,
\begin{equation}\label{eq:analytic-cubic-factorization}
 \zeta_K(s)=
  \prod_{d\in\{1,5,41,205\}}\prod_{j=0}^2
                L(s,\chi_d\vartheta^j)
  \prod_{j=0}^2 L(s,\psi\vartheta^j)^2.
\end{equation}
For clarity, the regular-representation Euler identity does
not discard bad primes: on inertia invariants, Frobenius acts
by cycles corresponding to the prime ideals over the given
rational prime, with their residue degrees. Its determinant
is therefore the product of the corresponding Dedekind Euler
factors. Direct-sum decomposition of that same invariant
space gives \eqref{eq:analytic-cubic-factorization} at every
prime.

There is no unproved Artin holomorphy input. Choose $C_4\subset
Q_8$ and a faithful linear character $\lambda$ of $C_4$.
Since its nontrivial conjugate is $\lambda^{-1}$,
$\operatorname{Ind}_{C_4}^{Q_8}\lambda=\psi$. Thus
\[
 \psi\vartheta^j=
 \operatorname{Ind}_{C_4\times C_3}^{Q_8\times C_3}
                       (\lambda\boxtimes\vartheta^j).
\]
Artin induction invariance identifies each $L$-function with
the nontrivial finite-order Hecke $L$-function over the
quadratic fixed field of $C_4\times C_3$. Such a function is
entire: the classical analytic input is, for example,
B.~Poonen, \emph{Tate's thesis}, section~5.11,
Theorem~5.22(b), p.~38, together with the Galois-character
construction in Example~5.14, p.~32. A nontrivial finite-order
character cannot be a norm power, the only exceptional
component in that theorem. The local induction identity is
the determinant identity for induced cosets, and holds at the
ramified places as well; it is the same unconditional
monomial argument already used for $\psi$ in NN-A018.

NN-A033 proves
$\operatorname{ord}_{1/2}L(s,\psi)=1$. Put
$k=\operatorname{ord}_{1/2}L(s,\psi\vartheta)\ge0$.
Conjugating the Artin Euler factors gives
\[
 L(s,\psi\bar\vartheta)
       =\overline{L(\bar s,\psi\vartheta)},
\]
because $\psi$ is real valued. This identity extends from
$\Re s>1$ to the entire continuations and shows that the
two cubic Artin twists have the same order $k$ at the center.
The four nonzero values in Lemma~\ref{lem:analytic-cubic-four-values}
and their conjugates dispose of the eight nonreal
degree-one factors. The remaining factors are $\zeta$ and
the three positive quadratic central values already certified
in NN-A018; $\zeta(1/2)<0$. Hence
\eqref{eq:analytic-cubic-factorization} gives exactly
\[
                  m=2+4k.
\]
It also determines the leading sign. On the real axis the
nonreal character factors occur as absolute squares. If
$b_1$ and $c_k$ are the first nonzero coefficients of
$L(s,\psi)$ and $L(s,\psi\vartheta)$, respectively, their
contribution to the leading coefficient is
$b_1^2|c_k|^4>0$. Multiplying by the negative nonzero product
of the degree-one central values gives
$a_m=\zeta_K^{(m)}(1/2)/m!<0$.

Finally the root discriminant is
\[
 D_K^{1/24}=205^{3/4}751^{2/3}>4371.
\]
This comparison is an exact integer inequality: raising to
the twelfth power gives
\[
 205^9 751^8-4371^{12}
    =16062670137855473202693706441001475807620284>0.
\]
The conductor criterion proved and interval-certified in
NN-A018 therefore gives $C_K>0$. Applying its full second
Puiseux formula with $a_m<0$ and $m=2+4k$ yields
\[
 \#\{\Re s_j(t)<1/2\}=\frac{m+2}2=2+2k,
 \quad
 \#\{\Re s_j(t)>1/2\}=\frac{m-2}2=2k
\]
on a common sufficiently small interval $t>0$. The two
vertical coefficients are $\pm iA$, and their second normal
coefficient is $-C_KA^2/m<0$. This proves all the assertions.
\end{proof}

\paragraph{Scope and arithmetic consequence.}
Together NN-A033 and NN-A034 supply explicit number fields
with opposite central side imbalances under the same
ideal-norm shift. The latter field is fixed before the
finite certification and requires no nonabelian GRH. The
unknown $k$ affects the total number of branches but not the
excess of two on the left. No noncentral sign-density law,
Riemann-zeta simplicity theorem or heat-collision statistic
is inferred from these central examples.

\paragraph{AG and representation translation.}
Away from $5\cdot41\cdot751$, the specified compositum gives
an actual finite etale $Q_8\times C_3$-cover of
$\operatorname{Spec}\mathbb Z[1/(5\cdot41\cdot751)]$.
The external tensor-product decomposition of its regular
representation gives the displayed Artin factorization.
Complex conjugation pairs the two cubic degree-two twists,
so their contribution to the central analytic length is
$4k$, while the old symplectic factor contributes two.
After $t=u^{2+4k}$, the local analytic Cartier divisor has
$2+4k$ normalized sections, with the two vertical first jets
having strictly negative second real normal component.
These are actual representation multiplicities and local
Artin lengths. The exterior-power/representation-ring
formalism is SGA~6, Expos\'e~V, Example~2.2(b),
Lemma~2.2.1; the local finite-cover/discriminant dictionary
is SGA~1, Expos\'e~I, section~4, Proposition~10. The
specified analytic norm shift is not asserted to arise
from an arithmetic deformation of the cover.
